\chapter{Fase 2: Prueba SIP con softphone}

\section{Objetivo}

Validar la centralita Asterisk antes de tocar el HT801. Para ello se registra un \emph{softphone} (Linphone) como extensión 102, se llama al número 9000 y se confirma que los tonos DTMF marcados desde la app llegan a Asterisk en forma de eventos AMI.

Este paso es importante porque permite verificar la parte de software del sistema sin depender aún del hardware analógico, simplificando el diagnóstico de problemas.

\section{Paso 1: Instalar Linphone}

Linphone es un cliente SIP gratuito y multiplataforma:

\begin{itemize}
  \item \textbf{iOS:} App Store, buscar \texttt{Linphone}.
  \item \textbf{Android:} Google Play, buscar \texttt{Linphone}.
  \item \textbf{macOS:} \texttt{linphone.org/technical-corner/linphone}.
\end{itemize}

Instalar la aplicación en cualquiera de los dispositivos disponibles que esté en la misma red local que la Pi.

\section{Paso 2: Crear cuenta SIP en Linphone}

Abrir Linphone y entrar al asistente: \textsf{Assistant} $\to$ \textsf{Use a SIP account} (o \textsf{Add Account}, según versión).

Completar:

\begin{table}[h]
\centering
\begin{tabular}{ll}
\toprule
\textbf{Campo} & \textbf{Valor} \\
\midrule
Username & 102 \\
Password & el secreto definido para la extensión 102 en \texttt{sip.conf} \\
Domain & 192.168.1.10 (IP de la Pi) \\
Transport & UDP \\
Display Name & Test Hugo \\
\bottomrule
\end{tabular}
\end{table}

Confirmar y verificar que el indicador de estado se ponga en verde con la leyenda \texttt{Registered} o \texttt{Connected}.

\section{Paso 3: Verificar registro desde Asterisk}

En la Pi:

\begin{shellcode}
sudo asterisk -rvvv
\end{shellcode}

Dentro del CLI:

\begin{plaincode}
hugo-pbx*CLI> sip show peers
Name/username             Host                Status
101/101                   (Unspecified)       UNKNOWN
102/102                   192.168.1.20:xxxx   OK (3 ms)
\end{plaincode}

La extensión 102 debe aparecer con la IP del dispositivo donde corre Linphone y estado \texttt{OK}.

\section{Paso 4: Llamada de prueba al 9000}

Desde Linphone, marcar \texttt{9000} y presionar el botón de llamada.

Se deben observar dos cosas simultáneamente:

\begin{enumerate}
  \item En el auricular del dispositivo: un \emph{beep} corto, producto del paso \texttt{Playback(beep)} en el dialplan.
  \item En la consola de Asterisk: la traza de ejecución del dialplan.
\end{enumerate}

\begin{plaincode}
-- Executing [9000@hugo-game:1] NoOp("PJSIP/102-...", "Hugo Phone — ...")
-- Executing [9000@hugo-game:2] Answer("PJSIP/102-...", "")
-- Executing [9000@hugo-game:3] Wait("PJSIP/102-...", "1")
-- Executing [9000@hugo-game:4] Playback("PJSIP/102-...", "beep")
-- <PJSIP/102-...> Playing 'beep.gsm' (language 'es')
-- Executing [9000@hugo-game:5] Read("PJSIP/102-...", "TECLA,,1,,,3600")
\end{plaincode}

\section{Paso 5: Probar DTMF}

Mientras la llamada sigue activa, en Linphone abrir el teclado de tonos y marcar el dígito \textbf{2}.

En la consola de Asterisk se debe ver:

\begin{plaincode}
== Manager 'hugo-bridge' logged on from 127.0.0.1
== UserEvent: HugoKey
   Channel: PJSIP/102-00000001
   Caller: 102
   Key: 2
-- Executing [9000@hugo-game:6] GotoIf(...)
-- Executing [9000@hugo-game:5] Read("PJSIP/102-...", "TECLA,,1,,,3600")
\end{plaincode}

Repetir con \textbf{4}, \textbf{6}, \textbf{8} y \textbf{5}. Cada uno debe disparar su propio \texttt{UserEvent}.

Si los eventos aparecen, la \textbf{Fase 2 está completa}: la centralita captura DTMF correctamente y el bridge podrá leerlos.

\section{Resolución de problemas}

\subsection{Linphone no se registra}

\begin{itemize}
  \item Verificar que la Pi sea alcanzable desde el dispositivo: \texttt{ping 192.168.1.10}.
  \item Confirmar que el puerto UDP 5060 no esté bloqueado por firewall en la Pi (Raspberry Pi OS Lite no trae firewall activado por defecto).
  \item Revisar logs de Asterisk en busca de intentos fallidos:
  \begin{shellcode}
sudo asterisk -rvvv
hugo-pbx*CLI> sip set debug on
  \end{shellcode}
\end{itemize}

\subsection{La llamada al 9000 se corta inmediatamente}

\begin{itemize}
  \item Verificar que la extensión 102 esté asignada al contexto \texttt{hugo-game} en \texttt{sip.conf}.
  \item Ejecutar \texttt{dialplan show hugo-game} en el CLI y confirmar que la extensión 9000 está presente.
\end{itemize}

\subsection{Conecta pero no se escucha el beep}

\begin{itemize}
  \item Problema típico de NAT o de codecs incompatibles. Verificar que \texttt{allow=ulaw} y \texttt{allow=alaw} estén en \texttt{sip.conf} y que Linphone también soporte estos codecs en su configuración avanzada.
  \item Revisar el volumen del dispositivo.
\end{itemize}

\subsection{No aparecen UserEvent al marcar DTMF}

\begin{itemize}
  \item Linphone puede enviar DTMF como tonos \emph{inband} en lugar de RFC 2833. En \textsf{Settings} $\to$ \textsf{Audio} $\to$ \textsf{DTMF mode} forzar el método \texttt{RFC 2833} o \texttt{SIP INFO}.
  \item Confirmar que el archivo \texttt{sip.conf} tiene la línea \texttt{dtmfmode=rfc2833} para la extensión 102.
\end{itemize}

\section{Cierre de fase}

En este punto está validado:

\begin{itemize}
  \item Asterisk acepta clientes SIP y los registra correctamente.
  \item El dialplan responde llamadas y reproduce audio.
  \item La función \texttt{Read()} captura DTMF en variable.
  \item El \texttt{UserEvent} se emite y es visible en AMI.
\end{itemize}

Esto es todo lo que el bridge necesita del lado de Asterisk. La siguiente fase prepara DOSBox-Staging y Hugo para que el bridge tenga algo donde inyectar las teclas.
